\begin{table}[t]
\centering
\caption{Held-out $R^2$ against two floor baselines, identical IDC and COAD train/held-out splits as Table~\ref{tab:validity} (seed=42). \texttt{Training-mean} predicts every held-out patch with the training set's per-gene mean (no image, no metadata). \texttt{Total-counts-linear} predicts from a single feature -- the patch's own raw (pre-normalization) total panel counts -- exactly the confound the library-size normalization fix targets; because normalization makes every patch's target sum to the same value by construction, this baseline can only detect a confound in \emph{relative} composition, not overall scale, so its small $R^2$ is a limited, not fully independent, check.}
\label{tab:baseline_comparison}
\begin{tabular}{llcc}
\toprule
Organ & Model & $R^2$ (whole panel) & $R^2$ (Hallmark genes) \\
\midrule
IDC & CONCH-Ridge (ours) & 0.339 & 0.402 \\
IDC & UNI-Ridge (ours) & 0.391 & 0.463 \\
IDC & Training-mean (no image) & -0.002 & -0.001 \\
IDC & Total-counts-linear (1 feature) & 0.068 & 0.070 \\
COAD & CONCH-Ridge (ours) & 0.400 & 0.495 \\
COAD & UNI-Ridge (ours) & 0.459 & 0.562 \\
COAD & Training-mean (no image) & -0.002 & -0.002 \\
COAD & Total-counts-linear (1 feature) & 0.080 & 0.122 \\
\bottomrule
\end{tabular}
\end{table}
